\documentclass{article}

\usepackage{amsmath}
\usepackage{graphicx}
\usepackage{float}
\title{Flood Interval Analysis}
\author{EA030}
\date{\today}

\begin{document}

\maketitle

\section*{Flood Interval (Return Period)}

In hydrology, the \textbf{flood interval} (or \textbf{return period}) represents the average time between floods of a certain magnitude. The return period \( T \) is defined as:

\[
T = \frac{n + 1}{m}
\]

Where:
\begin{itemize}
  \item \( T \) is the return period in years
  \item \( n \) is the number of years of record
  \item \( m \) is the rank of the flood event (1 = largest, 2 = second largest, etc.)
\end{itemize}

\section*{R Code: Flood Frequency Analysis}

Below is an R code chunk that demonstrates how to calculate flood intervals using annual peak discharge data.

\begin{knitrout}
\definecolor{shadecolor}{rgb}{0.969, 0.969, 0.969}\color{fgcolor}\begin{kframe}
\begin{alltt}
\hlcom{# Sample annual peak discharge data (in cubic meters per second)}
\hlstd{discharge} \hlkwb{<-} \hlkwd{c}\hlstd{(}\hlnum{400}\hlstd{,} \hlnum{520}\hlstd{,} \hlnum{610}\hlstd{,} \hlnum{700}\hlstd{,} \hlnum{650}\hlstd{,} \hlnum{720}\hlstd{,} \hlnum{800}\hlstd{,} \hlnum{750}\hlstd{,} \hlnum{680}\hlstd{,} \hlnum{670}\hlstd{)}

\hlcom{# Sort in descending order}
\hlstd{sorted_discharge} \hlkwb{<-} \hlkwd{sort}\hlstd{(discharge,} \hlkwc{decreasing} \hlstd{=} \hlnum{TRUE}\hlstd{)}

\hlcom{# Rank and return period}
\hlstd{rank} \hlkwb{<-} \hlnum{1}\hlopt{:}\hlkwd{length}\hlstd{(sorted_discharge)}
\hlstd{n} \hlkwb{<-} \hlkwd{length}\hlstd{(discharge)}
\hlstd{return_period} \hlkwb{<-} \hlstd{(n} \hlopt{+} \hlnum{1}\hlstd{)} \hlopt{/} \hlstd{rank}

\hlcom{# Create a data frame}
\hlstd{flood_data} \hlkwb{<-} \hlkwd{data.frame}\hlstd{(}
  \hlkwc{Rank} \hlstd{= rank,}
  \hlkwc{Discharge} \hlstd{= sorted_discharge,}
  \hlkwc{Return_Period} \hlstd{= return_period}
\hlstd{)}

\hlcom{# Print table}
\hlkwd{print}\hlstd{(flood_data)}
\end{alltt}
\end{kframe}
\end{knitrout}



\section*{Flood Data Table}

Below is the table of calculated return periods for each ranked discharge.





